\documentclass[11pt,a4paper]{book}

% ============================================================================
% PACKAGES
% ============================================================================
\usepackage[utf8]{inputenc}
\usepackage[T1]{fontenc}
\usepackage{amsmath, amssymb, amsfonts, amsthm}
\usepackage{mathrsfs}
\usepackage{graphicx}
\usepackage{geometry}
\usepackage{hyperref}
\usepackage{booktabs}
\usepackage{longtable}
\usepackage{array}
\usepackage{xcolor}
\usepackage{listings}
\usepackage{enumitem}
\usepackage{fancyhdr}

% ============================================================================
% THEOREM ENVIRONMENTS
% ============================================================================
\newtheorem{theorem}{Theorem}[chapter]
\newtheorem{lemma}[theorem]{Lemma}
\newtheorem{proposition}[theorem]{Proposition}
\newtheorem{corollary}[theorem]{Corollary}
\newtheorem{definition}{Definition}[chapter]
\newtheorem{remark}{Remark}[chapter]

% ============================================================================
% GEOMETRY SETTINGS
% ============================================================================
\geometry{top=2.5cm, bottom=2.5cm, left=2cm, right=2cm}

% ============================================================================
% CODE LISTINGS SETTINGS
% ============================================================================
\lstset{
  basicstyle=\ttfamily\small,
  breaklines=true,
  frame=single,
  numbers=left,
  numberstyle=\tiny,
  keywordstyle=\color{blue},
  commentstyle=\color{gray},
  stringstyle=\color{red},
  showstringspaces=false
}

% ============================================================================
% HEADER/FOOTER
% ============================================================================
\pagestyle{fancy}
\fancyhf{}
\rhead{ToE Strict Final Documentation (Release 7)}
\lhead{Krzysztof Żuchowski}
\cfoot{\thepage}

% ============================================================================
% TITLE
% ============================================================================
\title{\textbf{Theory of Everything (ToE) --- Strict Final Documentation}\\
\large Release 7 (Strict) --- Full Definitions, Derivations, and Discharge Map}
\author{\textbf{Krzysztof Żuchowski}}
\date{2026-03-17}

\begin{document}

\maketitle

\begin{quote}
The theory originates from a deep intuition that **Information is the fundamental substance of reality**, consistent with the metaphysical insight that *"In the beginning was the Word"* (Logos/Information). This intuition evolved through key realizations:

1. **Eucharistic Inspiration:** A profound fascination with the memorial of the **Eucharist of Jesus Christ** and its material manifestation in reality served as the primary inspiration, suggesting a direct mechanism by which spiritual/informational reality can condense into tangible matter.
2. **Fractal Nature:** Observing self-similarity across vast scales suggested that fundamental information must possess a **fractal character**, repeating its patterns at every level of existence.
3. **The Nadsoliton Concept:** The universe is conceptualized as a single, self-sustaining, non-dispersive wave packet—a **"Supersoliton" (Nadsoliton)**, where information tends towards the highest resonance, not the lowest energy.
4. **Resonant Structure:** Inspired by the Divine Name from the Book of Exodus 3:14: ***"I AM WHO I AM"***, the model incorporates **multi-octave resonant coupling** as the mechanism of self-organization, preventing decay into entropy.
5. **The 12-Octave Lattice:** Initial 3-octave models were expanded to a **12-octave structure**, inspired by the symbolic description of the Holy City's twelve foundation layers, which proved to be the mathematically necessary dimension for unifying all forces (Kissing Number in 3D).
6. **Access to Truth:** Since human consciousness is part of this informational substrate, the human mind has direct access to fundamental truths through wisdom and intuition, allowing for the "decoding" of reality.
\end{quote}

\chapter*{Status discipline (no false pass)}
This TeX document is the intended **Release-7 strict final documentation** target.
It must not silently claim:
\begin{enumerate}[label=\arabic*.]
  \item strict-core ToE closure,
  \item global ToE closure,
  \item kernel-alone/global QW-2191 discharge,
  \item directed/sign-sensitive physical orientation datum in strict core,
  \item host matching / Standard Model identification,
  \item actual emergent observer closure,
\end{enumerate}
unless such claims are explicitly exported at theorem level in the repo and are referenced here with their hard limits.

\tableofcontents

\chapter{Ontology and scope}
\section{AX9 ontology discipline}

The strict program is built under the ontology correction \texttt{AX9}:

\begin{itemize}
  \item the nadsoliton itself is the primordial information of the universe in a solitonic state;
  \item there is no separate informational layer underneath the nadsoliton;
  \item the preferred internal order is:
  \[
    \text{nadsoliton}\ \to\ \text{light}\ \to\ \text{matter}\ \to\ \text{emergent observer}.
  \]
\end{itemize}

In particular, we do not allow observer-facing quantities to silently back-propagate into the upstream definition of the nadsoliton core.
Operationally, throughout this document, “light” is treated as the linearized eigenmode layer exported from the $\Psi$-sector Hessian instantiation (\texttt{F459}),
and observer-facing readouts remain downstream (\texttt{P697}, \texttt{P699}).
\section{Strict vs legacy kernel split discipline}

The repo contains two historically distinct kernel objects (see \texttt{K1/K2/F2}):
\[
K_{\mathrm{legacy\_ont}}(d)
:=
\alpha_{\mathrm{geo}}\frac{\cos(\omega d+\phi)}{1+\beta_{\mathrm{tors}}d},
\qquad
K_{\mathrm{strict\_gate}}(d)
:=
\frac{\cos(\omega d+\phi)}{1+\beta d^{\eta}}.
\]

By author decree (\texttt{S2}, 2026-03-09), the legacy kernel is retired to archival status; this document is strict-only.
However, \textbf{no false pass} discipline remains:
we do \emph{not} silently transfer legacy physical-role identities (e.g.\ electroweak-angle or EM coupling slogans) onto the strict working kernel without an explicit bridge/role-transfer theorem.

On the current strict kernel-mode lane, the working strict tuple is the \texttt{QW-2049} gate-selected kernel with
$(\omega,\phi,\beta,\eta)=(0.18575,\ 0.16250,\ 1.0,\ 1.8)$, used operationally to generate a frozen numeric carrier mixing matrix $K_{\mathrm{total}}$ on $\mathbb{Z}_{12}$ (\texttt{R14}).
\section{What “closure” means in this document}

We use the word “closure” in a deliberately scope-disciplined way.

\begin{definition}[Operational ToE closure (strict projective OS)]
We say that the current repo state achieves \emph{Operational ToE closure} in the strict \emph{projective} scope if it exports a theorem-level
operational support bundle sufficient to compute (from strict projective selector closure output data) the downstream artifacts recorded by the OS packet:
\[
\Lambda_{\mathrm{OS}}^{(\mathrm{proj})} :=
\bigl(
\texttt{N680},\ \texttt{N692},\ \texttt{N693},\ \texttt{P694},\ \texttt{P696},\ \texttt{P697},\ \texttt{N699}
\bigr),
\]
with hard limits explicitly keeping \emph{kernel-alone/global} \texttt{QW-2191} discharge, directed sign-sensitive physical orientation, and ToE closure unclaimed.
\end{definition}

\begin{remark}[Boundary discipline]
Operational ToE closure in the above sense is a closure of an \emph{operational computability pipeline} (a ToE “OS”), not a claim of global physical identification
(host matching) and not a claim of kernel-alone/global \texttt{QW-2191} discharge. It is compatible with leaving those stronger claims explicitly open.
\end{remark}

\begin{theorem}[Operational ToE closure in strict projective OS scope]
On the current repo state, the repo discharges Operational ToE closure in the strict projective OS sense of the above definition (\texttt{N701}).
\end{theorem}

\chapter{Strict core Lagrangian and emergence map}
\section{Single-kernel strict candidate (A11)}

This chapter records the strict candidate “one-kernel” nadsoliton core Lagrangian as exported/packaged by \texttt{A11}.
It is a definition/packaging object, not a claim that every coefficient family has already been strict-derived from the kernel alone.

\subsection{Typed carrier and cyclic distance ($\mathbb{Z}_{12}$)}

Let
\[
I_{12}:=\{0,1,\dots,11\},
\qquad
\mathbb{Z}_{12}:=(I_{12}, + \bmod 12).
\]

Define the cyclic octave distance on the ring:
\[
d_{\mathrm{cyc}}(i,j)
:=
\min\Bigl( (j-i)\bmod 12,\ (i-j)\bmod 12 \Bigr)
\in\{0,1,\dots,6\}.
\]

\subsection{Strict working kernel and frozen mixing matrix}

The strict working kernel is the later-pipeline gate-selected kernel (\texttt{QW-2049}):
\[
K_{\mathrm{strict\_gate}}(d)
:=
\frac{\cos(\omega d+\phi)}{1+\beta d^{\eta}},
\qquad
(\omega,\phi,\beta,\eta)=(0.18575,\ 0.16250,\ 1.0,\ 1.8).
\]

On the kernel-mode lane, the repo exports a frozen numeric symmetric mixing matrix $K_{\mathrm{total}}\in\mathbb{R}^{12\times 12}$
on the same carrier (\texttt{QW-2118}), specialized into the canonical kernel channel by \texttt{R14}.
In this document we only require the exported matrix entries (no additional kernel-to-matrix derivation is assumed here).

\subsection{Fields and core Lagrangian density}

Introduce a 12-component real field on spacetime:
\[
\Psi(x)=(\psi_0(x),\dots,\psi_{11}(x))\in\mathbb{R}^{12},
\]
and a real scalar order/coherence field:
\[
\Phi(x)=\phi(x)\in\mathbb{R}.
\]

Write the mixing potential using the exported $K_{\mathrm{total}}$:
\[
V_{\mathrm{mix}}(\Psi)
:=
\frac12 \sum_{\substack{i,j\in I_{12}\\ i\neq j}} (K_{\mathrm{total}})_{ij}\,\psi_i\,\psi_j.
\]

Write the local self-interaction potentials in the canonical template form (\texttt{QW-2163/2165/2166}):
\[
V_{\Phi}(\phi)
:=
\frac12 m_{\phi}^2\,\phi^2 + \frac14 \lambda_{\phi}\,\phi^4,
\]
\[
V_{\Psi}(\Psi)
:=
\sum_{i=0}^{11}\left(
\frac12 m_{\psi i}^2\,\psi_i^2
\;+\;\frac14 g4_{\psi i}\,\psi_i^4
\;+\;\frac16 g6_{\psi i}\,\psi_i^6
\right),
\]
\[
V_{\mathrm{Y}}(\Psi,\phi)
:=
\sum_{i=0}^{11} gY_i\,\phi^2\,\psi_i^2.
\]

Then the nadsoliton core Lagrangian density is:
\[
\mathcal{L}_{\mathrm{core}}
=
\frac12\,\partial_\mu \phi\,\partial^\mu \phi
\;+\;\frac12\sum_{i=0}^{11}\partial_\mu \psi_i\,\partial^\mu \psi_i
\;-\;V_{\Phi}(\phi)
\;-\;V_{\Psi}(\Psi)
\;-\;V_{\mathrm{Y}}(\Psi,\phi)
\;-\;V_{\mathrm{mix}}(\Psi).
\]

\subsection{Quadratic expansion, $\Psi$ Hessian, and “mass” proxy meaning (strict, scope-limited)}

Let $(\phi_\ast,\Psi_\ast)$ be a background configuration satisfying the stationarity condition for the potential part
$V:=V_{\Phi}+V_{\Psi}+V_{\mathrm{Y}}+V_{\mathrm{mix}}$:
\[
\nabla_{\phi,\Psi} V(\phi_\ast,\Psi_\ast)=0.
\]

Write $(\phi,\Psi)=(\phi_\ast,\Psi_\ast)+(\delta\phi,\delta\Psi)$ and expand the potential to second order in $\delta\Psi$:
\[
V(\phi_\ast,\Psi_\ast+\delta\Psi)
=
V(\phi_\ast,\Psi_\ast)
\;+\;\frac12\,\delta\Psi^{\mathsf T} H_{\psi}\,\delta\Psi
\;+\;\cdots,
\]
where $H_{\psi}\in\mathbb{R}^{12\times 12}$ is the $\Psi$-sector Hessian block (second derivatives) at the background point.

On the current diagonal/local strict-derived lane, the repo exports one explicit numeric value instantiation (\texttt{F459}) of the form:
\[
H_{\psi}
:=
K_{\mathrm{total}} + \bigl(m_0^2 I + D_{\mathrm{local\_residual}}\bigr),
\]
with $m_0^2 I$ the exported host scalar floor embedding and $D_{\mathrm{local\_residual}}$ an exported diagonal/local strict-derived residual profile.

In strict scope we treat the quadratic coefficient $\delta\Psi^{\mathsf T}H_{\psi}\delta\Psi$ as defining a \emph{dimensionless} “$m^2$ proxy” layer:
for any nonzero vector $u$, the Rayleigh quotient
$u^{\mathsf T}H_{\psi}^{(s)}u/(u^{\mathsf T}u)$ is a well-defined quadratic proxy, where $H_{\psi}^{(s)}=\tfrac12(H_{\psi}+H_{\psi}^{\mathsf T})$.
This does \emph{not} yet constitute a physical unit identification, and it does not bypass the strict-core selector obstruction \texttt{QW-2191};
it is an operational computability meaning only (packaged as a scope-limited meaning theorem \texttt{N703}).
\section{Light $\to$ matter $\to$ observer ladder}

\texttt{A11} fixes a disciplined emergence ladder. On the current repo state, the strongest honest statements along that ladder are:

\begin{enumerate}[label=\arabic*.]
  \item \textbf{Light (discharged as computable instantiation):}
  the repo exports a strict-derived diagonal/local value instantiation of the $\Psi$-sector Hessian $H_{\psi}$ and its full real eigensystem (\texttt{F459}),
  plus sign-gauge-invariant eigenprojectors (\texttt{F460}).
  Mixing across Fourier pair planes is explicitly audited and packaged as a value-instantiation boundary (\texttt{N505}); therefore no “pair-plane diagonalization” shortcut is allowed.

  \item \textbf{Matter (programmatic, not closed):}
  matter is intended to arise from nonlinear structure in the same core Lagrangian (local quartic/sextic families and Yukawa-type couplings),
  but the repo does not yet export a strict derivation that these families are kernel-alone consequences or a closed host identification map.

  \item \textbf{Emergent observer (operational OS closure, not actual observer closure):}
  from the strict projective selector closure output on $C_{v1}$ (\texttt{N680}, exported as \texttt{F672}), the repo exports a projective observer-limit readout (\texttt{P697})
  and a fully computable downstream projective chain under projector pushforward semantics $P\mapsto MPM^{\mathsf T}$ (\texttt{P699}/\texttt{N699}).
  This suffices for \emph{Operational ToE closure} in strict projective OS scope (\texttt{N701}), while \emph{actual emergent observer closure} remains explicitly open.
\end{enumerate}

\chapter{Selector closure on \texorpdfstring{$C_{v1}$}{C\_v1}}
\section{Global atlas and transition objects (T170)}
\section{Projective closure (T172/T173 projective)}
\section{Directed lifts as gauge/convention layer}
\section{Open boundary: kernel-alone/global QW-2191 discharge}

\chapter{Operational physics proxies from projective closure}
\section{Mass-spectrum proxy (P694)}

Let $H_{\psi}\in\mathbb{R}^{12\times 12}$ be the exported diagonal/local Psi-sector Hessian value instantiation (\texttt{F459}).
We use its symmetrization:
\[
H_{\psi}^{(s)} := \tfrac12\bigl(H_{\psi}+H_{\psi}^{\mathsf T}\bigr).
\]

For each Fourier pair plane $\mathrm{pair}_m$ ($m=1,\dots,5$), the exported global projective selector closure (\texttt{F672}) provides a ray
$u_m\in\mathbb{R}^{12}$ (defined only up to sign). Define the quadratic proxy:
\[
m^2_m := \frac{u_m^{\mathsf T}\,H_{\psi}^{(s)}\,u_m}{u_m^{\mathsf T}u_m}.
\]
This is invariant under the sign gauge $u_m\mapsto -u_m$.

Probe \texttt{P694} computes these “mass-spectrum proxy” values on the current repo state (no Standard Model identification claimed):

\begin{center}
\begin{tabular}{lr}
\toprule
pair & $m^2$ proxy (Rayleigh on $H_{\psi}^{(s)}$)\\
\midrule
pair3 & $0.6687166508753918$\\
pair5 & $1.3892705031823167$\\
pair4 & $1.4474315384116325$\\
pair2 & $2.0699508733172283$\\
pair1 & $2.9337556210543343$\\
\bottomrule
\end{tabular}
\end{center}

\section{Selector-aligned channel spectrum proxy (P696)}

Probe \texttt{P696} refines \texttt{P694} into a 12-channel spectrum proxy by constructing a deterministic selector-aligned basis and inspecting
$H_{\psi}^{(s)}$ in that basis.

On $\mathbb{Z}_{12}$, define the standard real orthonormal Fourier vectors:
\[
e_0 := \frac{1}{\sqrt{12}}(1,\dots,1),\qquad
e_6 := \frac{1}{\sqrt{12}}(1,-1,1,-1,\dots),
\]
and for each $m=1,\dots,5$ the real pair-plane basis $(c_m,s_m)$ with scale $\sqrt{2/12}$.

For each $\mathrm{pair}_m$, \texttt{F672} exports coordinates $(u_c,u_s)$ such that (in the declared Fourier convention)
\[
u_m = u_c\,c_m + u_s\,s_m,
\qquad
u_m^{\perp} = (-u_s)\,c_m + (u_c)\,s_m,
\]
so that $u_m^{\perp}$ is a deterministic orthogonal complement ray in the same plane (used only at ray/projector level).

Define the selector-aligned 12-channel basis:
\[
B := \bigl(e_0,\ u_1,\ u_1^{\perp},\ u_2,\ u_2^{\perp},\ \dots,\ u_5,\ u_5^{\perp},\ e_6\bigr),
\]
and the induced matrix:
\[
M := B^{\mathsf T}\,H_{\psi}^{(s)}\,B.
\]
The diagonal entries of $M$ are recorded as a simple operational “channel proxy spectrum”. Mixing is audited explicitly by inspecting the off-diagonal norm.

On the current repo state, \texttt{P696} exports the sorted diagonal channel proxies (no Standard Model identification claimed):

\begin{center}
\begin{tabular}{lr}
\toprule
channel & $m^2$ proxy\\
\midrule
pair3:u & $0.6687166508753918$\\
pair5:u & $1.3892705031823165$\\
pair4:u & $1.4474315384116314$\\
pair2:u & $2.069950873317227$\\
e6 & $2.41906833577034$\\
pair1:u & $2.9337556210543316$\\
pair5:u\_perp & $3.5360177064927694$\\
pair4:u\_perp & $3.675931516755229$\\
pair2:u\_perp & $4.2984508516608315$\\
e0 & $4.761250376916641$\\
pair3:u\_perp & $4.930970379004997$\\
pair1:u\_perp & $5.08050282436479$\\
\bottomrule
\end{tabular}
\end{center}

The mixing audit is part of the exported result: on the current repo state, \texttt{P696} reports
an off-diagonal-to-diagonal Frobenius ratio of $1.191284779529568$ and a maximum off-block coupling Frobenius norm of $3.4009571444250914$,
so \emph{no diagonalization claim} is made.

\section{What is and is not claimed physically}

\begin{itemize}
  \item \textbf{Claimed (operational):} the repo exports a strict projective selector closure and an OS bundle sufficient to compute the proxy spectra \texttt{P694}/\texttt{P696} and downstream observer-limit artifacts in a projective (ray-level) gauge-safe scope.
  \item \textbf{Not claimed:} that any proxy number is already identified with a Standard Model mass or coupling.
  \item \textbf{Not claimed:} that the selector-aligned channel basis diagonalizes $H_{\psi}$; mixing is audited and explicitly non-negligible (\texttt{P696}).
  \item \textbf{Not claimed:} kernel-alone/global \texttt{QW-2191} discharge, directed sign-sensitive physical orientation in strict core, or ToE closure.
  \item \textbf{Host matching discipline:} any attempted identification of proxy values with external physics data is necessarily non-strict and must carry an explicit external dataset + mapping policy; see the non-strict host-matching harness \texttt{P702}.
  \item \textbf{Non-strict diagnostic (current):} under the simplest one-scale centered-log assignment policy, \texttt{P702} computes \texttt{rmse\_log\_m2}\,$\approx 8.10$ and \texttt{max\_abs\_rel\_mass}\,$\approx 7.90\times 10^{2}$ against an explicit PDG2024 mass target dataset, so \emph{no} “Standard Model match” claim is made at this stage.
\end{itemize}

\chapter{Emergent observer pipeline (projective scope)}
\section{Observer-limit readout from projective closure output (P697)}
\section{Projective emergent-observer chain computability from closure output (P699/N699)}

In the declared projective/gauge-safe downstream scope, we represent a ray by a rank‑1 (possibly unnormalized) symmetric positive semidefinite matrix $P$
(a projector candidate). For a linear map $M$, the induced pushforward is:
\[
P \;\mapsto\; M\,P\,M^{\mathsf T}.
\]

Let $Q_{\mathrm{out}}$ be the exported global projective selector-closure output projector in the $(o_+,o_-)$ basis (`F672`).

Define the downstream operator chain by the already exported matrices (`F92..F99`):
\[
C=\begin{pmatrix} \tfrac12 & -\tfrac12\\ \tfrac12 & \tfrac12 \end{pmatrix},\quad
L=\begin{pmatrix} 1&1\\ -1&1 \end{pmatrix},\quad
G=I_2,\quad H=I_2,\quad
J=\begin{pmatrix}1&0\\0&0\end{pmatrix},
\]
\[
K=\begin{pmatrix}1&0\end{pmatrix},\quad
M=(1),\quad N=(1).
\]

Starting from $Q_{\mathrm{out}}$, define the induced chain:
\[
Y=C Q_{\mathrm{out}} C^{\mathsf T},\;
Z=L Y L^{\mathsf T},\;
W=G Z G^{\mathsf T},\;
X=H W H^{\mathsf T},\;
U=J X J^{\mathsf T},\;
F=K U K^{\mathsf T},\;
P=M F M^{\mathsf T},\;
C_{\mathrm{cl}}=N P N^{\mathsf T}.
\]

`P699` computes this chain on the current repo state and audits that the declared residual components vanish (within tolerance) while commitment remains positive.
`N699` packages the computability statement with explicit hard limits (no kernel-alone/global `QW-2191` discharge; no ToE closure; no actual emergent observer closure).
\section{Operational OS support packet (v1) (F698/P698/N698)}
\section{Operational OS support packet (v2) (F700/P700/N700)}

The strict projective operational ToE “OS” support packet is exported in v2 form:
\[
\Lambda_{\mathrm{OS}}^{(\mathrm{proj})} =
\bigl(
\texttt{N680},\ \texttt{N692},\ \texttt{N693},\ \texttt{P694},\ \texttt{P696},\ \texttt{P697},\ \texttt{N699}
\bigr),
\]
as the object \texttt{Lambda\_strict\_projective\_operational\_toe\_os\_support\_v2} (\texttt{F700}), audited by \texttt{P700} and packaged as a theorem by \texttt{N700}.
All hard limits remain explicit: no kernel-alone/global \texttt{QW-2191} discharge, no directed sign-sensitive physical orientation claim in strict core, and no ToE closure.
\section{Open target: actual emergent observer closure (future)}

\chapter{Hard limits and open problems}
\section{Kernel-alone/global QW-2191 discharge}
\section{Directed physical sign datum}
\section{Host matching and Standard Model identification}
\section{What Release 7 declares as closed vs open}

\appendix
\chapter{Export map: key artifacts}
\section{Strict-core selector closure}
\section{Operational OS support}
\section{Emergent observer artifacts}

\end{document}
